\documentclass[a4paper,12pt]{article}
\usepackage[utf8]{inputenc}
\usepackage[spanish]{babel}
\usepackage{xcolor}
\usepackage{geometry}
\usepackage{tcolorbox}
\usepackage{titlesec}
\usepackage{graphicx}

% Estilo Ejecutivo
\definecolor{investorblue}{HTML}{1E3A8A}
\definecolor{successgreen}{HTML}{10B981}
\definecolor{lightbg}{HTML}{F1F5F9}

\geometry{margin=2.5cm}

\titleformat{\section}{\color{investorblue}\normalfont\Large\bfseries}{\thesection}{1em}{}

\begin{document}

\begin{titlepage}
    \centering
    \vspace*{2cm}
    {\Huge\bfseries\color{investorblue} Inteligencia Artificial de Próxima Generación}\\[0.5cm]
    {\Large\color{successgreen} Eficiencia Operativa y Ventaja Competitiva en Ganesha ERP}\\[2cm]
    
    \begin{tcolorbox}[colback=investorblue,colframe=investorblue,arc=8pt,left=20pt,right=20pt,top=20pt,bottom=20pt]
        \centering
        \color{white}
        \textbf{Propuesta de Valor Técnica para Inversores y Stakeholders}
    \end{tcolorbox}
    
    \vfill
    {\large \today}
\end{titlepage}

\section{El Desafío: ¿Por qué falla la IA tradicional?}
La mayoría de los sistemas de IA en la nube intentan "leer" toda la información de una empresa a la vez. Esto es costoso, lento y, a menudo, la IA se "confunde" con el volumen de datos (saturación).

\section{Nuestra Solución: "Inteligencia Directa"}
En Ganesha ERP, hemos rediseñado el cerebro del sistema (Backend) para que la IA no tenga que buscar a ciegas.

\begin{tcolorbox}[title=Innovación Clave: Canales de Datos Optimizados, colback=successgreen!5, colframe=successgreen]
    En lugar de darle a la IA miles de facturas sueltas, hemos creado \textbf{"resúmenes inteligentes"} automáticos. El servidor hace el trabajo pesado y le entrega a la IA la respuesta ya procesada.
\end{tcolorbox}

\section{Beneficios Estratégicos}
\begin{itemize}
    \item \textbf{Velocidad Extrema}: Las respuestas son casi instantáneas porque la IA no pierde tiempo analizando datos irrelevantes.
    \item \textbf{Reducción de Costos}: Al procesar los datos localmente en nuestro backend optimizado, no dependemos de costosos servicios externos de procesamiento por volumen.
    \item \textbf{Precisión del 100\%}: Al usar cálculos matemáticos directos del motor contable, eliminamos las "alucinaciones" típicas de la IA. Los números siempre cuadran.
\end{itemize}

\section{Escalabilidad y Futuro}
Nuestra arquitectura permite que el sistema crezca de 10 a 10,000 usuarios sin degradar la experiencia. Hemos optimizado la base de datos para que sea "amigable con la IA", permitiendo capturar tendencias y KPIs de forma fluida y sin saturar los recursos de la empresa cliente.

\section{Conclusión}
Ganesha ERP no es solo un software con un chat añadido; es una plataforma financiera diseñada desde su núcleo para ser operada por Inteligencia Artificial de forma eficiente, económica y segura.

\end{document}
